\section{Methods}

\subsection{Models}

TODO: Describe the two evaluated systems.

\begin{itemize}
  \item \textbf{Baseline model:} TODO, e.g. untouched Gemma model.
  \item \textbf{Ablated model:} TODO, e.g. same model after refusal-direction ablation.
\end{itemize}

The ablation procedure follows the idea that refusal behavior can be mediated by
a direction in activation space \citep{arditi2024refusal}. TODO: Briefly state
how the ablated model was produced, or state that the model outputs were provided
from a separate machine.

\subsection{Dataset}

The final analysis dataset contains one row per prompt. Each prompt is evaluated
on both models, giving paired observations. The binary response labels are:

\begin{equation}
  Y_{i,m} =
  \begin{cases}
    1, & \text{if model } m \text{ refused prompt } i,\\
    0, & \text{if model } m \text{ answered prompt } i.
  \end{cases}
\end{equation}

TODO: Insert final dataset counts in Table~\ref{tab:data-summary}.

\begin{table}[h]
  \centering
  \caption{Dataset summary. Replace the TODO values after the final labeled data is available.}
  \label{tab:data-summary}
  \begin{tabular}{lr}
    \toprule
    Quantity & Count \\
    \midrule
    Total prompts & TODO \\
    Unsafe prompts & TODO \\
    Benign prompts & TODO \\
    Prompt categories & TODO \\
    Median prompts per category & TODO \\
    \bottomrule
  \end{tabular}
\end{table}

\subsection{Metrics}

For unsafe prompts, a refusal is treated as the desired behavior and an answer is
treated as a non-refusal. Because the binary labels do not judge the content of
the answer, the unsafe-prompt answered rate is interpreted as a proxy for unsafe
compliance rather than a direct measure of harmfulness.

\begin{equation}
  \text{Unsafe non-refusal rate}_m =
  \frac{1}{N_{\text{unsafe}}}
  \sum_{i \in \text{unsafe}} (1 - Y_{i,m}).
\end{equation}

For benign prompts, refusal is treated as a false refusal:

\begin{equation}
  \text{False refusal rate}_m =
  \frac{1}{N_{\text{benign}}}
  \sum_{i \in \text{benign}} Y_{i,m}.
\end{equation}

We report rates for both models and the paired difference:

\begin{equation}
  \Delta = \bar{Y}_{\text{ablated}} - \bar{Y}_{\text{baseline}}.
\end{equation}

TODO: Specify whether confidence intervals are computed by bootstrap over
prompts or by an exact/binomial method for paired differences.

\subsection{Statistical Tests}

The primary statistical test is McNemar's test \citep{mcnemar1947note}. This is
appropriate because each prompt is evaluated on both models and the response
label is binary.

For unsafe prompts, the paired contingency table is:

\begin{table}[h]
  \centering
  \caption{Paired table for unsafe prompts.}
  \label{tab:mcnemar-unsafe-layout}
  \begin{tabular}{llcc}
    \toprule
    & & \multicolumn{2}{c}{Baseline model} \\
    \cmidrule(lr){3-4}
    & & Refused & Answered \\
    \midrule
    Ablated model & Refused & $a$ & $b$ \\
    Ablated model & Answered & $c$ & $d$ \\
    \bottomrule
  \end{tabular}
\end{table}

The null hypothesis is that the two off-diagonal flip probabilities are equal:

\begin{equation}
  H_0: P(\text{baseline refused, ablated answered})
  =
  P(\text{baseline answered, ablated refused}).
\end{equation}

For RQ1, the directional alternative is:

\begin{equation}
  H_A: P(\text{baseline refused, ablated answered})
  >
  P(\text{baseline answered, ablated refused}).
\end{equation}

For benign prompts, we use the same test structure, but the interpretation
changes: a refused benign prompt is counted as a false refusal.

\subsection{Category Analysis}

To answer RQ3, we report category-wise refusal-rate changes. If sample sizes are
sufficient, we also fit an exploratory logistic model:

\begin{equation}
  \text{logit}\left(P(Y_{i,m}=1)\right)
  =
  \beta_0 + \beta_1 \text{Ablated}_m +
  \beta_2 \text{Category}_i +
  \beta_3(\text{Ablated}_m \times \text{Category}_i).
\end{equation}

Because each prompt appears once for each model, uncertainty should account for
the paired design, for example by clustering standard errors by prompt id or by
using a repeated-measures model. The category analysis is treated as exploratory,
especially if many category-level comparisons are made.
